\par\Needspace{7\baselineskip}\bigskip
\begin{center}
\pdfbookmark[2]{Olbers to Gauss, Bremen, 5 April 1812}{pallas-brief-18120405}
\textsc{Olbers} to \textsc{Gauss}, Bremen, 5 April 1812.
\end{center}
\noindent\hspace*{\parindent}\{. . . . .\quad I am absolutely delighted by the exceedingly remarkable results you tell me of from your profound investigations of the orbit of Pallas. I too certainly consider your discovery that \(18\text{\AstroSymbols ♃}\) is exactly \(=7\text{\AstroSymbols ⚴}\) one of the most remarkable made in the solar system in a long time. — But what of Ceres? Does it have an appreciable equation depending on \(18\text{\AstroSymbols ♃}-7\text{\AstroSymbols ⚳}\)? Might similar equations of the form \(m\text{\AstroSymbols ♃}=n\text{\AstroSymbols ⚳}\), etc., be found for it and the other small planets as well?\quad . . . . .\}
\par\Needspace{7\baselineskip}\bigskip
\begin{center}
\pdfbookmark[2]{Gauss to Bessel, Göttingen, 5 May 1812}{pallas-brief-18120505}
\textsc{Gauss} to \textsc{Bessel}, Göttingen, 5 May 1812.
\end{center}
\noindent\hspace*{\parindent}. . . . .\ I have been occupied chiefly with the perturbations of Pallas by Jupiter. You will have read something about them in No. 67 of our Gelehrten Anzeigen. I am glad to tell you the remarkable result set down there in cipher, but ask that for the time being it remain \textls[180]{entirely between the two of us}. It is that the mean motions of {\AstroSymbols ♃} and {\AstroSymbols ⚴} stand in the \textls[180]{rational} ratio \(7:18\), which is continually restored \textls[180]{exactly} by Jupiter's action, as with the rotation period of our Moon. I have already begun a second calculation of the periodic perturbations.\quad . . . . .
\par\Needspace{7\baselineskip}\bigskip
\begin{center}
\pdfbookmark[2]{Olrers to Gauss, Paris, 12 May 1812}{pallas-brief-18120512}
\textsc{Olrers} to \textsc{Gauss}, Paris, 12 May 1812.
\end{center}
\noindent\hspace*{\parindent}\{. . . . .\quad \textsc{Laplace} asks in great detail about you and your work on Pallas. He seems very curious about your method, about which I could tell him nothing, nor did I wish to. You cannot imagine how enthusiastically he and all the others . . . . .\ welcomed the hope of persuading you . . . . .\ to accompany me on one of my future journeys to Paris.\quad . . . . .\}
\par
